\chapter{System Evaluation on the Evaluation Set}
\label{chap:system_evaluation}

\section{Experimental Results}
Before applying any score calibration or fusion strategies, it is crucial to benchmark the raw outputs of our selected optimal classifiers. We extracted the best-performing models from the previous laboratories to establish our baseline systems for the final project evaluation:
\begin{enumerate}
    \item \textbf{Quadratic Logistic Regression:} Prior-weighted Quadratic Logistic Regression ($\lambda=3.16 \times 10^{-2}$, $\tilde{\pi}=0.5$).
    \item \textbf{Support Vector Machine (SVM):} RBF Kernel SVM ($C=1.0$, $\gamma=0.01$).
    \item \textbf{Gaussian Mixture Model (GMM):} Tied Covariance GMM ($M=8$ components).
\end{enumerate}

These models were trained on the designated training split and evaluated on the held-out evaluation dataset. We assessed their performance across three distinct target applications: $\tilde{\pi} \in \{0.1, 0.5, 0.9\}$. The Minimum DCF and Actual DCF metrics for the uncalibrated (raw) scores are reported in Table \ref{tab:raw_eval_results}.

\begin{table}[H]
    \centering
    \caption{Baseline Performance (Raw Scores) on the Evaluation Set.}
    \label{tab:raw_eval_results}
    \resizebox{0.7\textwidth}{!}{
    \begin{tabular}{llcc}
        \toprule
        \textbf{Model} & \textbf{Application Prior ($\tilde{\pi}$)} & \textbf{minDCF} & \textbf{actDCF} \\
        \midrule
        \multirow{3}{*}{\textbf{Quadratic Logistic Regression}} 
        & $0.1$ & 0.3520 & 0.4920 \\
        & $0.5$ & 0.1410 & 0.1430 \\
        & $0.9$ & 0.3050 & 0.4910 \\
        \midrule
        \multirow{3}{*}{\textbf{RBF Kernel SVM}} 
        & $0.1$ & 0.4589 & 0.6550 \\
        & $0.5$ & 0.1877 & 0.1910 \\
        & $0.9$ & 0.4315 & 0.6751 \\
        \midrule
        \multirow{3}{*}{\textbf{Tied GMM (8 Comp.)}} 
        & $0.1$ & 0.4072 & 0.4257 \\
        & $0.5$ & \textbf{0.1656} & \textbf{0.1666} \\
        & $0.9$ & 0.3654 & 0.3727 \\
        \bottomrule
    \end{tabular}
    }
\end{table}

\section{Analysis and Observations}

\subsection*{Question: Are the models well calibrated for the given applications?}

Overall, the models exhibit varying and often severe degrees of miscalibration depending on the application prior. Calibration is measured by the gap between the \textit{Actual DCF} (the operational error using theoretical Bayes thresholds) and the \textit{Minimum DCF} (the lowest possible error achievable by searching for an optimal, data-driven threshold).

For the Quadratic Logistic Regression model, the calibration suffers significantly at extreme priors. At $\tilde{\pi}=0.1$, the actDCF ($0.4920$) is notably higher than the minDCF ($0.3520$), and similarly at $\tilde{\pi}=0.9$ (actDCF $0.4910$ vs minDCF $0.3050$). The SVM model suffers similarly at extreme priors, jumping from a minDCF of $0.4589$ to an actDCF of $0.6550$ at $\tilde{\pi}=0.1$. The generative GMM, however, shows significantly tighter bounds, demonstrating strong intrinsic calibration across all thresholds.

\subsection*{Question: Are there models that are better calibrated than others for the different applications?}

Yes, the \textbf{Tied GMM} is spectacularly better calibrated than the discriminative models (LR and SVM). For the GMM at $\tilde{\pi}=0.5$, the gap between minDCF ($0.1656$) and actDCF ($0.1666$) is a mere $0.0010$. Even at the heavily skewed $\tilde{\pi}=0.9$, the GMM's actDCF ($0.3727$) stays incredibly close to its minDCF ($0.3654$).

The \textbf{RBF Kernel SVM} is remarkably well calibrated \textit{only} at $\tilde{\pi}=0.5$ (gap of $0.0033$), but its calibration degrades catastrophically at extreme priors ($\tilde{\pi}=0.1$ and $\tilde{\pi}=0.9$). Finally, the \textbf{Quadratic Logistic Regression} model is very well calibrated at $\tilde{\pi}=0.5$, but its calibration diverges significantly at extreme priors, requiring proper score-to-LLR mapping to recover its true discriminative power.

\section{Theoretical Considerations}

The sharp differences in calibration behavior between these three classifiers can be directly explained by their underlying mathematical assumptions and the specific geometric properties of our dataset, which we have meticulously documented in previous laboratories.

\textbf{Generative vs. Discriminative Score Outputs:}
The Tied GMM is a generative model. By definition, it explicitly estimates the class-conditional probability densities $f_{X|C}(x|c)$. Because it models the true density of the data, the difference between the log-likelihoods of the two classes mathematically equates to a true Log-Likelihood Ratio (LLR). Consequently, comparing a GMM score to the theoretical Bayes threshold $-\log \frac{\tilde{\pi}}{1-\tilde{\pi}}$ is mathematically rigorous. As we discovered in Laboratory 4 and 5, our features exhibit severe non-Gaussianity and localized variance disparities. The multi-modal components of the GMM successfully capture these irregularities, resulting in highly accurate density estimation and, by extension, near-perfect native calibration.

\textbf{The Margin-based Nature of SVMs:}
Conversely, the Support Vector Machine is not a probabilistic model. It is a strictly discriminative, margin-maximizing algorithm. The "score" it produces is simply the geometric distance of a sample from the separating hyperplane in the kernel-induced feature space. The SVM has absolutely no concept of probability densities or priors. At $\tilde{\pi}=0.5$, the theoretical threshold is 0, which perfectly coincides with the SVM's hyperplane ($w^T \phi(x) = 0$), explaining why it appears well-calibrated at $\tilde{\pi}=0.5$. However, at $\tilde{\pi}=0.1$ or $0.9$, the theoretical Bayes threshold shifts non-linearly. Because SVM distances do not scale like true LLRs, comparing an SVM margin distance to a probabilistic threshold yields severe calibration loss.

\textbf{Quadratic Logistic Regression and $L_2$ Regularization:}
Unlike the strictly linear baseline, the Quadratic Logistic Regression expands the feature space to natively model non-linear separation boundaries, successfully capturing the misaligned covariance structures of our dataset. However, its scores are still heavily miscalibrated at extreme application priors. The true cause of this calibration loss is not an overly simplistic linear boundary, but rather the $L_2$ regularization term ($\lambda = 3.16 \times 10^{-2}$). The regularization heavily penalizes large weights during optimization to prevent overfitting. This artificial compression of the weight vector mathematically squashes the output scores toward zero, preventing them from growing large enough to reflect extreme confidence levels. Consequently, the output scores are structurally distorted and no longer function as well-calibrated log-likelihood ratios, directly leading to the observed calibration penalty.
